% BHU PYQ -- Professional Exam Template
\documentclass[11pt,a4paper]{article}

\usepackage[a4paper,top=2.5cm,bottom=2.5cm,left=2.8cm,right=2.8cm,headheight=14pt]{geometry}
\usepackage{amsmath,amssymb}
\usepackage[T1]{fontenc}
\usepackage[utf8]{inputenc}
\usepackage{lmodern}
\usepackage{microtype}
\usepackage{fancyhdr}
\usepackage{enumitem}
\usepackage{listings}
\usepackage{hyperref}
\usepackage{lastpage}
\usepackage{parskip}

\hypersetup{colorlinks=true,urlcolor=black,linkcolor=black,
  pdftitle={GEN-301: Digital Logic and Processors -- BHU PYQ}}

\pagestyle{fancy}
\fancyhf{}
\renewcommand{\headrulewidth}{0.4pt}
\renewcommand{\footrulewidth}{0.4pt}
\fancyhead[L]{\small\textit{Banaras Hindu University}}
\fancyhead[R]{\small\textit{GEN-301: Digital Logic and Processors}}
\fancyfoot[C]{\small Page \thepage\ of \pageref{LastPage}}

\lstset{
  basicstyle=\small\ttfamily,
  frame=single,
  breaklines=true,
  columns=flexible,
  keepspaces=true
}

\newlist{parts}{enumerate}{2}
\setlist[parts,1]{label=(\alph*),leftmargin=*,itemsep=4pt,topsep=3pt}
\setlist[parts,2]{label=(\roman*),leftmargin=*,itemsep=2pt,topsep=2pt}
\newcommand{\pts}[1]{\hfill\textit{[#1]}}

\begin{document}

\begin{center}
  {\large\bfseries BANARAS HINDU UNIVERSITY}\\[2pt]
  {\normalsize B.Voc. Semester III Examination, 2022-23}\\[6pt]
  \rule{0.8\linewidth}{0.5pt}\\[6pt]
  {\large\bfseries Paper: GEN-301 --- Digital Logic and Processors}\\[4pt]
  {\normalsize Time: Three hours \hfill Maximum Marks: 70}
\end{center}

\rule{\linewidth}{0.4pt}

\smallskip
\noindent\textit{Instructions: Attempt all questions. Symbols have their usual meanings.}

\bigskip

COMPUTER APPLICATION
GEN - 301 : Digital Logic and Processors
Question Paper)
> Answer All Sections.
> Candidates are required to write their answers in'their own words as far
as practicable.
. Section - A : (15x2=30) a
Long Answer type questions to be answered in about 750 words each :

\medskip\noindent\textbf{Question 1.} Prove the following Boolean identities :
\begin{parts}
  \item AC+ABC=AC
  \begin{parts}
    \item A+AB=A+B
    \item ABC+ABC+ABC —
    = A(BtC)
  \end{parts}
\end{parts}
\centerline{\textbf{OR}}

\medskip\noindent\textbf{Question 2.} Explain about the signal processing and types of it.

\medskip\noindent\textbf{Question 3.} Convert the number system from one form to another.
G@ (6)0=( pb
  \begin{parts}
    \item (110010), =  )s
    \item (A996 = ( dro
  \end{parts}
\centerline{\textbf{OR}}

\medskip\noindent\textbf{Question 4.} What are the basic gates in logic circuit. Give the brief explanation
. about them with truth table.

\medskip\noindent\textbf{Question 1.} em
\begin{parts}
  \item 
\end{parts}

\bigskip\noindent{\large\bfseries SECTION --- B : (5x6=30}
\rule{\linewidth}{0.4pt}\smallskip

Attempt any six:
Short answer type questions to be answered in 550 words if they are
theoretical.

\medskip\noindent\textbf{Question 1.} What is combinational circuit. Design any combinational! circuit with
two AND gates and one OR gate and draw the truth table.

\medskip\noindent\textbf{Question 2.} Perform mathematical operation : |
\begin{parts}
  \item 1101 (i) 1111
  +1010 - 101
\end{parts}

\medskip\noindent\textbf{Question 3.} What will be the output for the given logic circuit if X, = 1 and X, = 0.

\medskip\noindent\textbf{Question 4.} Calculate the number of cells in K-map if we have 3 pts and draw
the K-map as well as write the output in each cells.

\medskip\noindent\textbf{Question 5.} Demorganize the expression. |
(A+ B)(C+D)

\medskip\noindent\textbf{Question 6.} Convert the given logical expression in SOP.
ABC + ACD + ABCD

\medskip\noindent\textbf{Question 7.} What are the latches and what are it’s types ? Explain about NAND
latch.

\medskip\noindent\textbf{Question 8.} What is multiplexer ?
sen) ae    ((c))

\bigskip\noindent{\large\bfseries SECTION --- C : (10x1=10)}
\rule{\linewidth}{0.4pt}\smallskip


\medskip\noindent\textbf{Question 1.} In K-map what is the abbreviation of K ?

\medskip\noindent\textbf{Question 2.} Give the one example of universal gate.

\medskip\noindent\textbf{Question 3.} What type of arithmetic operation performed by OR gate ?

\medskip\noindent\textbf{Question 4.} The binary variable can store the value of zero or one. This statement
is true or false.

\medskip\noindent\textbf{Question 5.} Which device converts single input into the multiple output ?

\medskip\noindent\textbf{Question 6.} Does flip-flop store the 1-bit of information ?

\medskip\noindent\textbf{Question 7.} Does register use for storing the data ? €

\medskip\noindent\textbf{Question 8.} Multiply (1001), with (10).
° Q.9 Write the full form of MSB.

\medskip\noindent\textbf{Question 10.} Write the 1’s complement of (1000),.
ReREK

\medskip\noindent\textbf{Question 3.}
Ke

\vfill
\rule{\linewidth}{0.4pt}
\begin{center}\small
\href{https://bhu-exam.vercel.app/}{Click here to visit our website}\\[4pt]\href{https://me-aryan.vercel.app/}{Click here to visit our contributor}\\[4pt]\href{https://quashan-me.vercel.app/}{Click here to visit our contributor}
\end{center}

\end{document}
